\begin{abstract}
Editable, biologically interpretable 3D models of field-grown crops
support high-throughput phenotyping, \emph{in silico} breeding trials,
and agronomic digital twins, but recovering such models from field
sensor data remains difficult: plant organs are slender, neighboring
plants occlude one another, and a field campaign rarely yields more than
a single photograph, a noisy depth frame, or a partial point cloud per
plant.
% (fix 2: vision hook — delete the next sentence if unwanted)
Slender geometry under occlusion is a regime in which learned 3D
completion and Chamfer-based evaluation both fail. In this work, we
present \method{}, a multimodal
masked-autoencoder framework that recovers the complete procedural
representation of a plant, the generator parameters encoding leaf
count, insertion heights, blade angles, and surface geometry, from any
subset of RGB, depth, and point-cloud input in a single forward pass.
The model is pretrained on 15{,}000 procedurally generated plants per
species (maize and sorghum), fine-tuned on rows in which simulated
neighbors occlude the target, and distilled from a full-modality teacher
into single-sensor students.
Distillation lowers the cross-modal generation loss by 30\% in sorghum and
44\% in maize, and halves the sorghum parameter error from a point cloud alone.
Training on rows removes 92\% (maize) and 85\% (sorghum) of the error that
neighbors add for a model trained on isolated plants.
Controlled ablations in both species show that exposing most of the
procedural parameters to the encoder degrades phenotype recovery, while
using them only as a reconstruction target degrades the sorghum
representation and improves the maize one. We further show that Chamfer distance and linear probes on
the frozen representation rank input modalities differently on slender
geometry, and we adopt the probe as the deciding metric. On real field
scans with independently reconstructed ground truth, \method{}
\todo{matches/approaches} the accuracy of a staged annotation pipeline
and reduces the time required for procedural reconstruction from hours of manual
specification, or minutes of pipeline computation, to milliseconds per
plant. We will release \dataset{}, the \todo{N}-plant, two-species benchmark
behind these experiments, which pairs multimodal field captures under
simulated row occlusion with exact procedural ground truth.
\end{abstract}